% Author: Dr. Matthias Jung, DL9MJ
% Zweiton-/Eintonvergleich: NF, AM, SSB
%
% Benötigt:
% \usepackage{pgfplots}
% \usepgfplotslibrary{fillbetween}
% \pgfplotsset{compat=1.18}

\begin{tikzpicture}

    \pgfplotsset{
        samples=1000,
        every axis/.style={
            height=2.0cm,
            scale only axis,
            width=\linewidth*\getDarcImageFactor,
            axis lines=none,
            ticks=none,
            xmin=0,
            xmax=8.8,
            ymin=-2.1,
            ymax=2.1,
            domain=0:8.8,
            thick,
            smooth,
            no markers,
            title style={
                at={(0,1.02)},
                anchor=south west,
                font=\small
            }
        }
    }

    % normierte Frequenzen
    \pgfmathsetmacro{\fm}{0.70}   % Einton-NF, z.B. 700 Hz
    \pgfmathsetmacro{\fa}{0.70}   % erster Zweiton, z.B. 700 Hz
    \pgfmathsetmacro{\fb}{1.20}   % zweiter Zweiton, z.B. 1200 Hz
    \pgfmathsetmacro{\fc}{7.00}   % HF-Träger, nur für Darstellung
    \pgfmathsetmacro{\m}{0.70}    % AM-Modulationsgrad

    % ------------------------------------------------------------
    % 1. NF-Eintonsignal
    % ------------------------------------------------------------
    \begin{axis}[
        name=plot1,
        title={1. Nutzsignal NF: Einton},
        at={(0,0)}
    ]
        \addplot[DARCred]
            {sin(deg(2*pi*\fm*x))};
    \end{axis}

    % ------------------------------------------------------------
    % 2. AM von 1
    % ------------------------------------------------------------
    \begin{axis}[
        name=plot2,
        title={2. AM von 1: Einton moduliert Trägeramplitude},
        at={(0,-2.8cm)}
    ]
        \addplot[name path=amoneupper, draw=none]
            {1 + \m*sin(deg(2*pi*\fm*x))};

        \addplot[name path=amonelower, draw=none]
            {-(1 + \m*sin(deg(2*pi*\fm*x)))};

        \addplot[gray, opacity=0.35]
            fill between[of=amoneupper and amonelower];

        \addplot[thick, black]
            {1 + \m*sin(deg(2*pi*\fm*x))};

        \addplot[thick, black]
            {-(1 + \m*sin(deg(2*pi*\fm*x)))};
    \end{axis}

    % ------------------------------------------------------------
    % 3. SSB von 1
    % ------------------------------------------------------------
    \begin{axis}[
        name=plot3,
        title={3. SSB von 1: eine HF-Linie, konstante Hüllkurve},
        at={(0,-5.6cm)}
    ]
        \addplot[name path=ssboneupper, draw=none] {1};
        \addplot[name path=ssbonelower, draw=none] {-1};

        \addplot[gray, opacity=0.18]
            fill between[of=ssboneupper and ssbonelower];

        \addplot[thick, black] {1};
        \addplot[thick, black] {-1};
    \end{axis}

    % ------------------------------------------------------------
    % 4. Zweiton-Nutzsignal
    % ------------------------------------------------------------
    \begin{axis}[
        name=plot4,
        title={4. Zweiton-Nutzsignal NF: z.B. 700 Hz und 1200 Hz},
        at={(0,-8.4cm)}
    ]
        \addplot[DARCred]
            {0.5*(sin(deg(2*pi*\fa*x))
                + sin(deg(2*pi*\fb*x)))};
    \end{axis}

    % ------------------------------------------------------------
    % 5. AM von 4
    % ------------------------------------------------------------
    \begin{axis}[
        name=plot5,
        title={5. AM von 4: Zweiton moduliert Trägeramplitude},
        at={(0,-11.2cm)}
    ]
        \addplot[name path=amtwoupper, draw=none]
            {1 + \m*0.5*(sin(deg(2*pi*\fa*x))
                        + sin(deg(2*pi*\fb*x)))};

        \addplot[name path=amtwolower, draw=none]
            {-(1 + \m*0.5*(sin(deg(2*pi*\fa*x))
                         + sin(deg(2*pi*\fb*x))))};

        \addplot[gray, opacity=0.35]
            fill between[of=amtwoupper and amtwolower];

        \addplot[thick, black]
            {1 + \m*0.5*(sin(deg(2*pi*\fa*x))
                        + sin(deg(2*pi*\fb*x)))};

        \addplot[thick, black]
            {-(1 + \m*0.5*(sin(deg(2*pi*\fa*x))
                         + sin(deg(2*pi*\fb*x))))};
    \end{axis}

    % ------------------------------------------------------------
    % 6. SSB von 4
    % ------------------------------------------------------------
    \begin{axis}[
        name=plot6,
        title={6. SSB von 4: zwei HF-Linien mit Schwebung},
        at={(0,-14.0cm)}
    ]
        \addplot[name path=ssbtwoupper, draw=none]
            {abs(cos(deg(pi*(\fb-\fa)*x)))};

        \addplot[name path=ssbtwolower, draw=none]
            {-abs(cos(deg(pi*(\fb-\fa)*x)))};

        \addplot[gray, opacity=0.35]
            fill between[of=ssbtwoupper and ssbtwolower];

        \addplot[thick, black]
            {abs(cos(deg(pi*(\fb-\fa)*x)))};

        \addplot[thick, black]
            {-abs(cos(deg(pi*(\fb-\fa)*x)))};
    \end{axis}
\end{tikzpicture}